\documentclass[a4paper,10pt]{article}
\usepackage[utf8]{inputenc}
\usepackage{amsmath, amssymb}
\usepackage[top=1cm, bottom=1cm, left=1cm, right=1cm]{geometry}
\renewcommand{\arraystretch}{1.2}

\begin{document}
\pagestyle{empty}

% Primera evaluación
\begin{minipage}[t]{0.45\textwidth}
    \begin{center}

        \textbf{Nombre y Apellido: \underline{\hspace{4cm}}}
    \end{center}


    \begin{enumerate}
        \item \textbf{(2 pts)} Un objeto de 500 g de masa se sumerge completamente en agua. Calcule la fuerza de flotación que actúa sobre el objeto. (Densidad del agua = 1000 kg/m³)
        
        

            \item \textbf{(2 pts)} El caudal de un río es de 500 m³/s. Si el ancho del río es de 50 m y la profundidad es de 2 m, ¿cuál es la velocidad del agua en el río?
        
        \item \textbf{(2 pts)} Una manguera de jardín tiene un diámetro de 2 cm y se estrecha a un diámetro de 1 cm. Si el agua fluye a 1 m/s en la sección más ancha, ¿cuál es la velocidad del agua en la sección estrecha?
        
        \item \textbf{(2 pts)} Un cilindro de 20 cm de altura y 5 cm de diámetro flota en agua con su eje vertical. Si la densidad del cilindro es de 600 kg/m³, ¿qué altura del cilindro está sumergida?
        
  
        
        \item \textbf{(2 pts)} Un iceberg flota en agua de mar con una densidad de 1025 kg/m³. Si la densidad del hielo es de 917 kg/m³, ¿qué fracción del volumen del iceberg está por encima del nivel del agua?
    \end{enumerate}
\end{minipage}
\hfill
% Segunda evaluación
\begin{minipage}[t]{0.45\textwidth}
    \begin{center}

        \textbf{Nombre y Apellido: \underline{\hspace{4cm}}}
    \end{center}

    \begin{enumerate}
        \item \textbf{(2 pts)} Un tubo de Venturi tiene un diámetro de entrada de 10 cm y un diámetro de estrechamiento de 5 cm. Si la velocidad del agua en la entrada es de 2 m/s, ¿cuál es la velocidad del agua en el estrechamiento?

        \item \textbf{(2 pts)} Un cubo de 10 cm de lado está completamente sumergido en agua. Determine el empuje que experimenta.

      \item \textbf{(2 pts)} Un barco de 2000 kg flota en agua dulce. Calcule el volumen de agua desplazado por el barco. (Densidad del agua = 1000 kg/m³)
        \item \textbf{(2 pts)} Un tanque de agua tiene una pequeña fuga en el fondo. Si el tanque tiene 2 m de altura y está lleno, ¿con qué velocidad sale el agua por la fuga? (Suponga que la velocidad del agua en la superficie del tanque es despreciable)
        
        \item \textbf{(2 pts)} Un fluido incompresible fluye a través de una tubería horizontal que se estrecha de un diámetro de 12 cm a un diámetro de 6 cm. Si la velocidad del fluido en la sección más ancha es de 3 m/s, determine la velocidad del fluido en la sección estrecha.
        
    
    \end{enumerate}
\end{minipage}




\end{document}
